% Part: normal-modal-logic
% Chapter: syntax-and-semantics
% Section: relational-models

\documentclass[../../../include/open-logic-section]{subfiles}

\begin{document}

\olfileid{nml}{syn}{rel}

\olsection{关系模型}

正规模态逻辑语义的基本概念是\emph{关系模型}。
它由一个世界集合、一个将这些世界联系起来的二元“可达关系”，
以及一个确定哪些命题变元在哪些世界上“为真”的赋值构成。

\begin{defn}
  基本模态语言的一个\emph{模型}是一个三元组 $\mModel{M}
  = \tuple{W, R, V}$，其中
  \begin{enumerate}
  \item $W$ 是一个非空的“世界”集合，
  \item $R$ 是 $W$ 上的二元可达关系，
  \item $V$ 是一个函数，它为每个命题变元~$p$ 指派一个可能世界的集合 $V(p)$。
  \end{enumerate}
  当 $Rww'$ 成立时，我们称 $w'$ \emph{从 $w$ 可达}。
  当 $w \in V(p)$ 时，我们称 $p$ \emph{在 $w$ 为真}。
\end{defn}

关系语义的一大优势在于，模型可以用简单的图表来表示，例如图~\olref{fig:simple}。
世界由节点表示，且世界~$w'$ 从~$w$ 可达当且仅当存在一条从~$w$ 指向~$w'$ 的箭头。
此外，当 $w \in V(p)$ 时，我们用~$\mTrue{p}$ 标记该节点（世界），否则用~\mFalse{p} 标记。
图~\olref{fig:simple} 表示了具有 $W = \{w_1, w_2, w_3\}$、$R = \{\tuple{w_1, w_2},\tuple{w_1,
  w_3}\}$、$V(p) = \{w_1, w_2\}$ 和 $V(q) = \{w_2\}$ 的模型。

\begin{figure}
  \begin{center}
    \begin{tikzpicture}[modal]
      \node[world] (w1) [label={[align=right]right:\mTrue{p}\\ \mFalse{q}}]{$w_1$};
      \node[world] (w2) [label={[align=right]right:\mTrue{p}\\ \mTrue{q}},
        above right=of w1]{$w_2$};
      \node[world] (w3) [label={[align=right]right:\mFalse{p}\\ \mFalse{q}},
        below right=of w1] {$w_3$};
      \draw[->] (w1) to (w2);
      \draw[->] (w1) to (w3);
    \end{tikzpicture}
  \end{center}
  \caption{一个简单模型。}
  \ollabel{fig:simple}
\end{figure}

\end{document}